\section{Motorsteuerung} \label{cpt:Motorsteuerung}


\subsection{Grundlagen zur Motorsteuerung}


\subsection{TwinCAT Motion Baukasten}

TwinCAT Motion ist eine Komponente des TwinCAT-Automatisierungssystems von Beckhoff, das speziell für die Steuerung von Bewegungsachsen und Antrieben in Maschinen und Produktionsanlagen entwickelt wurde. Die Tc2\_Mc2-Bibliothek ist ein integraler Bestandteil der TwinCAT Motion-Umgebung und bietet Funktionen für die Steuerung und Regelung von Achsen und Bewegungen. Hier werden die angeschlossenen Achsen und Antriebe konfiguriert. Dies umfasst die Definition von Achsentypen, Regelungsparametern, Verfahrbereichen und anderen relevanten Einstellungen. Es wird eine Auswahl verschiedener Regelungsalgorithmen für die Achsensteuerung, darunter P-PI-D-Regler  und erweiterte Regelungen wie Zustandsregelungen zur Verfügung gestellt. Zusätzlich sind umfassende Diagnose- und Überwachungsfunktionen, einschließlich Echtzeit-Visualisierung von Achspositionen, Geschwindigkeiten, Beschleunigungen und Fehlerstatus möglich. \cite{InfosysNCMotion}

Die Tc2\_Mc2-Bibliothek bietet leistungsstarke Funktionen zur Bewegungssteuerung. Dies umfasst die Definition von Bewegungsprofilen, Geschwindigkeitsprofilen, Beschleunigungsprofilen und Positionierungsparametern. Die Motion-Steuerung ist eng in die TwinCAT-PLC-Programmierung integriert. Programmierer können auf die Tc2\_Mc2-Bibliothek zugreifen und die Motion-Funktionalität direkt in ihrer IEC-61131-3-konformen Programmiersprache nutzen. \cite{InfosysTc2Mc2}



\subsubsection{Reglerimplementierung}

\subsection{Basisaufgaben}

\subsubsection{Förderband - 5 Punkte}

\textbf{DC-Motor}


\textbf{Inkrementalencoder}


\textbf{NC-Parameter Konfiguration}

\subsubsection{Extruder - 5 Punkte}

\textbf{Schrittmotor}

\textbf{NC-Parameter Konfiguration}

\subsubsection{Manipulatorachsen - 5 Punkte}

\newpage
